\chapter{Conclusion}
\label{ch:conclusion}

This thesis asked whether longitudinal transitions in the fuzzy cluster membership of football players have incremental explanatory power for changes in their market values, using the robust fuzzy $C$-medoids model for mixed data of \citeA{durso2022robust} on eight seasons of the Portuguese Primeira Liga. The analysis combined FBref statistics for 3,192 player-seasons of 1,499 outfield players with Transfermarkt valuations for 94.5\% of the player-seasons, clustered every season separately, aligned the archetypes across seasons, and related changes in memberships to changes in log market value for 1,361 pairs of consecutive seasons, with a bootstrap that accounts for the estimation of the memberships. The objective was explanatory and associational, and no causal claim is made.

The clustering identifies four stable and interpretable archetypes, defenders, central midfielders, wide attackers and strikers, and a small noise cluster. Players are fairly persistent in their archetype (75\% remain), with defenders almost fixed and attacking players more mobile. The estimation of the attribute weights, a distinctive feature of the model, failed on these data and had to be replaced by weights fixed in proportion to the number of variables in each group, and the alignment of the archetypes across seasons had to use their profiles instead of their medoids.

Against the two pre-specified baselines, the membership transitions are weakly associated with changes in market value in the main configuration: the bootstrap $p$-values are 0.015 against the full baseline and 0.079 against the core baseline, and the gains are small (0.3 to 0.6 percentage points of adjusted $R^2$ and 0.1 to 0.3 points of out-of-sample $R^2$). The association is concentrated in the shift from the defender profile to the wide-attacker profile. It is also fragile. It is found much less often with five archetypes than with four, it varies by orders of magnitude between random draws of the clustering, a typical draw brings no improvement in out-of-sample prediction over the core baseline, and no result survives a correction for the number of specifications examined. Above all, it largely disappears when the changes in the position labels, which enter the clustering but not the baselines, are controlled: the bootstrap $p$-values are then 0.35 and 0.12, and the position changes alone add more explanatory power than the memberships do. The conclusion is therefore that membership transitions carry no robust incremental explanatory power for changes in market value, and that what looks like such power in the main run largely reflects changes in the position labels.

The thesis contributes a longitudinal application of a mixed-data fuzzy clustering model to a league outside the big five, a documented solution to the failure of weight estimation and to the inconsistency of medoid-based alignment, a bootstrap that accounts for estimated memberships, and a full report of how a seemingly positive association varies across settings, random draws and controls. Its limitations are the restricted set of statistics available, the noise in the memberships and in market values, the limited sample, and a specification search in which an important control was added only after an independent review. For practitioners, the archetypes are a useful description but the transitions are not a basis for valuation. For research, replication in larger leagues with the tuning and the controls fixed in advance and the clustering repeated over many seeds, a test of the league-relative position of a player as the one source of genuinely new information, and methods that treat the noise in memberships directly, are the natural next steps.
